\documentclass[12pt,a4paper]{article} % Use A4 paper with a 12pt font size - different paper sizes will require manual recalculation of page margins and border positions 

\usepackage{marginnote} % Required for margin notes
\usepackage{wallpaper} % Required to set each page to have a background
\usepackage{lastpage} % Required to print the total number of pages
\usepackage[left=1.3cm,right=2.0cm,top=1.8cm,bottom=5.0cm,marginparwidth=3.4cm]{geometry} % Adjust page margins
\usepackage{amsmath} % Required for equation customization
\usepackage{amssymb} % Required to include mathematical symbols
\usepackage{xcolor} % Required to specify colors by name

\usepackage{fancyhdr} % Required to customize headers
\usepackage[brazil]{babel}
% \usepackage[latin1]{inputenc}
%\usepackage[T1]{fontenc} 
\usepackage{graphicx}
\usepackage{pstricks}
\usepackage{subfigure}
\usepackage{caption}  % legendas nas figuras
\captionsetup{justification=centering,labelfont=bf}
\usepackage{textcomp}
\setlength{\headheight}{80pt} % Increase the size of the header to accommodate meta-information
\pagestyle{fancy}\fancyhf{} % Use the custom header specified below
\renewcommand{\headrulewidth}{0pt} % Remove the default horizontal rule under the header

\setlength{\parindent}{0cm} % Remove paragraph indentation
\newcommand{\tab}{\hspace*{2em}} % Defines a new command for some horizontal space

\newcommand\BackgroundStructure{ % Command to specify the background of each page
\setlength{\unitlength}{1mm} % Set the unit length to millimeters

\setlength\fboxsep{0mm} % Adjusts the distance between the frameboxes and the borderlines
\setlength\fboxrule{0.5mm} % Increase the thickness of the border line
\put(10, 20pr){\fcolorbox{black}{gray!5}{\framebox(155,247){}}} % Main content box
\put(165, 20){\fcolorbox{black}{gray!10}{\framebox(37,247){}}} % Margin box
\put(10, 262){\fcolorbox{black}{white!10}{\framebox(192, 25){}}} % Header box
\put(175, 263){\includegraphics[height=23mm,keepaspectratio]{}} % Logo box - maximum height/width: 
}

%----------------------------------------------------------------------------------------
%	HEADER INFORMATION
%----------------------------------------------------------------------------------------

\fancyhead[L]{\begin{tabular}{l r | l r} % The header is a table with 4 columns
\textbf{MI} & TEC501 & \textbf{P\'agina:} & \thepage/\pageref{LastPage} \\ % Project name and page count
\textbf{Atividade:} & Relat\'orio & \textbf{Data:} & 18/09/13 \\ % Job number and last updated date
\textbf{Professor:} & Delmar B. Carvalho & \textbf{Assunto:} & Resposta ao Impulso  \\ % Version and reviewed date
\textbf{Prazo Entrega:} & \textbf{03/10/13} & \textbf{Autor(es):} & Aluno 1; Aluno2 \\ % Designer and reviewer
\end{tabular}}

%----------------------------------------------------------------------------------------

\begin{document}

%\AddToShipoutPicture{\BackgroundStructure} % Set the background of each page to that specified above in the header information section

%----------------------------------------------------------------------------------------
%	DOCUMENT CONTENT
%----------------------------------------------------------------------------------------

\section{Relat\'orio t\'ecnico descritivo sobre ......} 

Este modelo reduzido visa proporcionar um roteiro pr\'atico para a reda\,c\~ao e apresentação dos resultados obtidos de experimentos, desenvolvimentos ou
equacionamentos, na forma de relat\'orio simplificado. Esta se\,c\~ao \'e destinada a apresentar os objetivos do trabalho.\\

\subsection{T\'itulo da sub-se\,c\~ao}

No desenvolvimento da atividade, caso alguma ideia precise ser desdobrada poder\~ao ser criadas sub-se\,c\~oes descritivas do processo. Tamb\'em poder\~ao ser inseridas
equa\, c\~os em linha como: $s=\jmath\omega$ ou ent\~ao atrav\'es de um bloco de equa\,c\~ao como este sem numera\,c\~ao autom\'atica:
\begin{equation*} % o * serve para bloquear a numeracao automatica
\tag{x.1} % destinada a designar uma sequencia ou informacaoo complementar da equacao
T(s)=\frac{Y(s)}{X(s)}
\label{eq:funcaoTransferencia}
\end{equation*}

ou ent\~ao com numera\,c\~ao autom\'atica:

\begin{equation} 
H(s)=\frac{1}{T(s)}
\label{eq:funcaoAtenucao}
\end{equation}
\\

As equa\,c\~oes apresentadas podem ser citadas no texto apartir dos respectivos indexadores criados com o comando 
\begin{verbatim}
\label
\end{verbatim}
inserido na equa\,c\~ao. Assim a equa\,c\~ao que representa a atenua\,c\~ao, Eq.(\ref{eq:funcaoAtenucao}), ir\'a aparecer no texto, independente da sequ\^encia que ir\'a assumir. 
Da mesma forma a equa\,c\~ao da fun\,c\~ao de transefer\^encia (Eq. \ref{eq:funcaoTransferencia}) tamb\'em poder\'a ser citada.                                                                                                                                                                                                                                                                               

\par\vspace{\baselineskip}

%----------------------------------------------------------------------------------------

\end{document}